\section{Deep Learning for Sparse 3D Point Clouds}
\label{sec:deep-learning-3d}

\subsection{Sparse Convolutions}
\label{ssec:sparse-conv}

Standard 3D convolutions applied to a $G^3$ voxel grid have memory and
computational complexity $O(G^3)$, which is $O(10^9)$ for $G = 1024$.
Sparse convolutions~\cite{graham2015sparse,choy2019minkowski} restrict
computation to the active (occupied) voxels, achieving $O(N)$ complexity
where $N \ll G^3$ is the number of occupied voxels.

\paragraph{Generalised sparse convolution.}
For an input sparse tensor $(\mathbf{C}^{in}, \mathbf{F}^{in})$ and a kernel
$\mathbf{W}$ with offsets $\mathcal{O} = \{o_1, \ldots, o_{k^3}\}$
(the $k^3$ cells of a $k\times k\times k$ kernel), the output at coordinate
$u \in \mathcal{C}^{out}$ is:
\begin{equation}
  \mathbf{f}^{out}_u = \sum_{i : (u + o_i) \in \mathcal{C}^{in}}
    \mathbf{W}_i \mathbf{f}^{in}_{u + o_i} + \mathbf{b}.
  \label{eq:sparse-conv}
\end{equation}
The inner sum runs only over the subset of kernel offsets $o_i$ for which
the corresponding input coordinate $u + o_i$ is active.
The output coordinates $\mathcal{C}^{out}$ are determined by the kernel map
(the precomputed mapping from output to active input coordinates), which is
built once during the forward pass and reused for the backward pass.

\paragraph{Stride-2 sparse convolution (downsampling).}
A stride-2 sparse convolution halves the spatial resolution of the active
coordinate set: each output coordinate $u \in \mathcal{C}^{out}$ corresponds
to input coordinates in a $k\times k\times k$ neighbourhood of $2u$
(in the input coordinate system).
This reduces the number of active coordinates by approximately a factor of 8,
while doubling the effective receptive field.

\paragraph{Transposed sparse convolution (upsampling).}
A transposed (stride-2) sparse convolution inverts the downsampling:
given a coarse active coordinate set, it expands each coordinate to up to 8
output coordinates, recovering the original spatial resolution.
The output feature at each expanded coordinate accumulates contributions from
nearby input features according to the transposed kernel map.

\paragraph{Submanifold sparse convolution.}
A submanifold sparse convolution~\cite{graham20183d} preserves the active
coordinate set exactly: $\mathcal{C}^{out} = \mathcal{C}^{in}$.
This avoids the activations spreading out from the surface into empty space
during processing.
All intermediate (same-stride) convolutions in our backbone use submanifold
mode; only downsampling and upsampling layers change the active set.
Figure~\ref{fig:sparse-conv} contrasts dense, generalised sparse, and submanifold
convolutions on a 2D grid analogy.

\begin{figure}[t]
\centering
\resizebox{0.92\linewidth}{!}{%
\begin{tikzpicture}[>=Stealth, every node/.style={font=\scriptsize}]
  \def\cs{0.55}  % cell size

  % ===== (a) Dense convolution =====
  % Input: all cells filled
  \foreach \i in {0,...,4}
    \foreach \j in {0,...,4} {
      \fill[blue!18]  ({\j*\cs},{-\i*\cs}) rectangle ++({\cs},{\cs});
      \draw[gray!45, very thin] ({\j*\cs},{-\i*\cs}) rectangle ++({\cs},{\cs});
    }
  % RF: rows i=1..3, cols j=1..3 — fully inside grid
  \draw[red!80, very thick] ({\cs},{-3*\cs}) rectangle ++({3*\cs},{3*\cs});
  \node[below, align=center] at ({2.5*\cs},{-5*\cs-0.18})
    {(a) Dense input\\(all active)};

  % Arrow
  \node[font=\normalsize] at ({5*\cs+0.4},{-2.5*\cs}) {$\to$};

  % Dense output (all active)
  \begin{scope}[xshift={5*\cs cm + 0.75cm}]
    \foreach \i in {0,...,4}
      \foreach \j in {0,...,4} {
        \fill[green!20] ({\j*\cs},{-\i*\cs}) rectangle ++({\cs},{\cs});
        \draw[gray!40, very thin] ({\j*\cs},{-\i*\cs}) rectangle ++({\cs},{\cs});
      }
    \node[below, align=center] at ({2.5*\cs},{-5*\cs-0.18})
      {Dense output\\(all active)};
  \end{scope}

  % ===== (b) Sparse (generalised) convolution =====
  \def\spoff{6.80}  % x-offset for sparse section (in cm)

  % Separator
  \draw[gray!35, dashed, thin] ({\spoff-0.25},{ 0.6}) -- ({\spoff-0.25},{-3.5});

  % Sparse input
  \begin{scope}[xshift=\spoff cm]
    \foreach \i in {0,...,4}
      \foreach \j in {0,...,4}
        \draw[gray!28, very thin] ({\j*\cs},{-\i*\cs}) rectangle ++({\cs},{\cs});
    % Active input cells (sparse pattern)
    \foreach \i/\j in {0/2, 1/1, 1/4, 2/0, 2/3, 3/2, 3/4, 4/1} {
      \fill[blue!30] ({\j*\cs},{-\i*\cs}) rectangle ++({\cs},{\cs});
      \draw[blue!50, very thin] ({\j*\cs},{-\i*\cs}) rectangle ++({\cs},{\cs});
    }
    % RF: rows i=1..3, cols j=1..3 — same position as panel (a)
    \draw[red!80, very thick] ({\cs},{-3*\cs}) rectangle ++({3*\cs},{3*\cs});
    \node[below, align=center] at ({2.5*\cs},{-5*\cs-0.18})
      {(b) Sparse input\\(few active)};
  \end{scope}

  % Arrow
  \node[font=\normalsize] at ({\spoff+5*\cs+0.4},{-2.5*\cs}) {$\to$};

  % Sparse output (expanded active set)
  \begin{scope}[xshift={\spoff cm + 5*\cs cm + 0.75cm}]
    \foreach \i in {0,...,4}
      \foreach \j in {0,...,4}
        \draw[gray!28, very thin] ({\j*\cs},{-\i*\cs}) rectangle ++({\cs},{\cs});
    % Expanded output cells (more active than input)
    \foreach \i/\j in {0/1,0/2,0/3, 1/0,1/1,1/2,1/3, 2/0,2/1,2/2,2/3,2/4, 3/1,3/2,3/3,3/4} {
      \fill[green!32] ({\j*\cs},{-\i*\cs}) rectangle ++({\cs},{\cs});
      \draw[green!60, very thin] ({\j*\cs},{-\i*\cs}) rectangle ++({\cs},{\cs});
    }
    \node[below, align=center] at ({2.5*\cs},{-5*\cs-0.18})
      {Output expands\\(more active)};
  \end{scope}

\end{tikzpicture}%
}
\caption{%
  Dense convolution vs.\ sparse convolution (2D analogy).
  (a) Dense: kernel applied at every position, output is fully dense.
  (b) Generalised sparse: kernel applied only where input is active (blue),
  but output can activate new positions within each kernel's reach (shown in green).
  Submanifold sparse convolution (not shown) restricts $\mathcal{C}^{out} = \mathcal{C}^{in}$,
  preventing activation spread.
}
\label{fig:sparse-conv}
\end{figure}

\subsection{Residual Blocks for Sparse Tensors}

Residual blocks~\cite{he2016deep} are adapted straightforwardly to sparse
tensors by replacing dense 2D convolutions with submanifold sparse convolutions:
\begin{equation}
  \mathbf{F}^{out} = \mathbf{F}^{in} +
    \text{SubConv}_{3^3}\!\left(
      \text{ReLU}\!\left(
        \text{BN}\!\left(
          \text{SubConv}_{3^3}(\mathbf{F}^{in})
        \right)
      \right)
    \right),
\end{equation}
where BN denotes batch normalisation applied to the feature matrix $\mathbf{F}$
(across the $N$ active coordinates per channel) and SubConv$_{3^3}$ is a
submanifold $3\times3\times3$ sparse convolution.
The skip connection is exact: since SubConv preserves the active coordinate set,
the input and output feature matrices have the same number of rows and can
be added directly.
Figure~\ref{fig:resblock} shows the sparse residual block structure.

\begin{figure}[t]
\centering
\begin{tikzpicture}[>=Stealth, every node/.style={font=\small},
  box/.style={draw, rounded corners=3pt, inner sep=6pt, align=center,
              minimum width=3.6cm, minimum height=0.60cm}]
  % Main vertical path
  \node[box, fill=gray!10]  (in)  at (0,  0.0) {$\mathbf{F}^{in}$\enspace$(N \times C)$};
  \node[box, fill=blue!12]  (c1)  at (0, -1.3) {SubConv$_{3^3}$,\ BN,\ ReLU};
  \node[box, fill=blue!12]  (c2)  at (0, -2.6) {SubConv$_{3^3}$,\ BN};
  \node[circle, draw, fill=orange!22, inner sep=5pt] (add) at (0,-3.7) {$+$};
  \node[box, fill=gray!10]  (out) at (0, -5.0) {$\mathbf{F}^{out}$\enspace$(N \times C)$};
  % Main path arrows
  \draw[->] (in)  -- (c1);
  \draw[->] (c1)  -- (c2);
  \draw[->] (c2)  -- (add);
  \draw[->] (add) -- node[right, font=\scriptsize]{ReLU} (out);
  % Identity skip connection (routes cleanly to the right)
  \draw[->, blue!55, thick]
    (in.east) -- ++(1.8,0) -- ++(0,-3.7) -- (add.east);
  \node[rotate=-90, font=\scriptsize, blue!70] at (2.15, -1.85) {identity skip};
\end{tikzpicture}
\caption{%
  Sparse residual block for point cloud processing.
  Two submanifold $3^3$ sparse convolutions (with batch normalisation and ReLU)
  are stacked; the input is added via an identity skip connection.
  Because SubConv preserves the active coordinate set, the skip connection
  requires no projection or resampling.
}
\label{fig:resblock}
\end{figure}

\subsection{Batch Normalisation}
\label{ssec:batch-norm}

Batch normalisation~\cite{ioffe2015batch} normalises each feature channel
across the batch and spatial dimensions:
\begin{equation}
  \hat{f}^{(c)}_{n,i} = \frac{f^{(c)}_{n,i} - \mu^{(c)}}{\sqrt{(\sigma^{(c)})^2 + \epsilon}},
  \qquad
  \text{BN}(f^{(c)}_{n,i}) = \gamma^{(c)} \hat{f}^{(c)}_{n,i} + \beta^{(c)},
\end{equation}
where $\mu^{(c)}$ and $\sigma^{(c)}$ are the mean and standard deviation of
channel $c$ across all active coordinates in the mini-batch, and
$\gamma^{(c)}, \beta^{(c)}$ are learned affine parameters.
For sparse tensors, the batch statistics are computed across all $N \times B$
active coordinates in a mini-batch of size $B$, where $N$ varies per sample.
At inference time, running statistics accumulated during training are used.

\subsection{Encoder-Decoder Architecture (U-Net)}
\label{ssec:unet-theory}

The U-Net architecture~\cite{ronneberger2015unet} consists of a contracting
encoder path, a bottleneck, and an expanding decoder path with skip connections.
Applied to sparse point clouds with sparse convolutions, this gives:

\begin{align}
  \text{Encoder level } \ell: &\quad
    \mathbf{F}^{enc}_\ell = \text{SubConv}(\mathbf{F}^{enc}_{\ell-1}),
    \quad \mathbf{F}^{enc}_{\ell+1} = \text{StrideConv}(\mathbf{F}^{enc}_\ell) \\
  \text{Decoder level } \ell: &\quad
    \mathbf{F}^{dec}_\ell = \text{SubConv}([\mathbf{F}^{enc}_\ell;\, \text{Up}(\mathbf{F}^{dec}_{\ell+1})])
\end{align}
where $\text{Up}$ denotes transposed stride-2 convolution and $[\cdot\,;\,\cdot]$
is feature concatenation.

The key motivation for skip connections in point cloud processing is the same
as in image segmentation: the encoder's strided convolutions discard
high-resolution detail in exchange for larger receptive fields; skip connections
restore that detail to the decoder, which must produce high-resolution
per-point predictions.
Without skip connections, the network can only use low-resolution context
from the bottleneck, severely limiting its ability to localise fine geometric
corrections.
Figure~\ref{fig:unet-concept} illustrates the encoder-decoder pattern with
skip connections.

\begin{figure}[t]
\centering
\begin{tikzpicture}[>=Stealth, every node/.style={font=\small},
  box/.style={draw, rounded corners=2pt, minimum width=1.5cm,
              minimum height=0.5cm, inner sep=4pt, align=center}]
  % Encoder (left column)
  \node[box, fill=blue!15] (e1) at (0,  0)   {$\mathbf{F}_1$\\{\tiny 32 ch}};
  \node[box, fill=blue!20] (e2) at (0, -1.3) {$\mathbf{F}_2$\\{\tiny 64 ch}};
  \node[box, fill=blue!30] (e3) at (0, -2.6) {$\mathbf{F}_3$\\{\tiny 128 ch}};
  % Decoder (right column, mirrored)
  \node[box, fill=teal!30] (d3) at (4, -2.6) {$\mathbf{G}_3$\\{\tiny 128 ch}};
  \node[box, fill=teal!20] (d2) at (4, -1.3) {$\mathbf{G}_2$\\{\tiny 64 ch}};
  \node[box, fill=teal!15] (d1) at (4,  0)   {$\mathbf{G}_1$\\{\tiny 32 ch}};
  % Output
  \node[box, fill=gray!15] (out) at (4, 1.3) {Output};
  % Encoder downward path
  \draw[->] (e1) -- node[right,font=\scriptsize]{$\times 2\!\downarrow$} (e2);
  \draw[->] (e2) -- node[right,font=\scriptsize]{$\times 2\!\downarrow$} (e3);
  % Bottleneck (straight horizontal at bottom level)
  \draw[->, thick] (e3) -- node[above,font=\scriptsize]{bottleneck} (d3);
  % Decoder upward path
  \draw[->] (d3) -- node[right,font=\scriptsize]{$\div 2\!\uparrow$} (d2);
  \draw[->] (d2) -- node[right,font=\scriptsize]{$\div 2\!\uparrow$} (d1);
  \draw[->] (d1) -- (out);
  % Skip connections (horizontal at matching levels)
  \draw[->, dashed, blue!60, thick] (e1.east) --
    node[above,font=\scriptsize,blue!80]{skip} (d1.west);
  \draw[->, dashed, blue!60, thick] (e2.east) --
    node[above,font=\scriptsize,blue!80]{skip} (d2.west);
  % Column labels
  \node[blue!60, font=\small\bfseries, left] at (-1.0,-1.3) {Encoder};
  \node[teal!70!black, font=\small\bfseries, right] at (5.1,-1.3) {Decoder};
\end{tikzpicture}
\caption{%
  U-Net encoder-decoder with skip connections.
  The encoder contracts the spatial resolution while expanding the channel
  dimension (building global context); the decoder expands back to full
  resolution.
  Skip connections (dashed) concatenate high-resolution encoder features
  directly to the decoder, preserving fine-grained spatial detail that the
  bottleneck representation has lost.
}
\label{fig:unet-concept}
\end{figure}

\subsection{Feature-wise Linear Modulation (FiLM)}
\label{ssec:film}

Feature-wise Linear Modulation (FiLM)~\cite{perez2018film} is a general
mechanism for conditioning a neural network on an external context signal,
without modifying the core network architecture.
Given intermediate feature maps $\mathbf{h} \in \R^{N \times C}$ (where $N$
is the batch dimension and $C$ the channel dimension) and a conditioning input
$\mathbf{c}$ (a scalar, vector, or embedding), FiLM applies a learned
per-channel affine transformation:
\begin{equation}
  \text{FiLM}(\mathbf{h};\,\mathbf{c}) = \boldsymbol{\gamma}(\mathbf{c}) \odot \mathbf{h}
    + \boldsymbol{\beta}(\mathbf{c}),
  \qquad \boldsymbol{\gamma},\, \boldsymbol{\beta} \in \R^C,
\end{equation}
where $\boldsymbol{\gamma}(\mathbf{c})$ and $\boldsymbol{\beta}(\mathbf{c})$
are produced by a conditioning network (typically a small MLP) applied to $\mathbf{c}$.
The $\odot$ denotes channel-wise scaling (broadcast over spatial dimensions).

\paragraph{Historical context.}
FiLM was introduced~\cite{perez2018film} for visual question answering, where
the conditioning signal is a language embedding encoding the question.
The key insight was that allowing the question to modulate intermediate visual
feature maps (rather than only entering at the first or last layer),
dramatically improving the model's ability to selectively attend to question-relevant
image regions.
The mechanism has since been applied to multi-task
learning~\cite{rebuffi2018efficient}, style transfer, and physics
simulation~\cite{gupta2023pde}, in each case using a domain or parameter
embedding to steer a shared backbone without updating its weights.

In this work, the conditioning signal is the compression bitrate
($r = \log_2(\bpp)$), allowing a single post-processing model to adapt its
displacement predictions across multiple rate-distortion operating points
without separate models per rate.
Figure~\ref{fig:film-diagram} illustrates the FiLM conditioning mechanism.

\begin{figure}[t]
\centering
\begin{tikzpicture}[>=Stealth, every node/.style={font=\small},
  box/.style={draw, rounded corners=3pt, inner sep=5pt, align=center,
              minimum height=0.58cm},
  op/.style={circle, draw, inner sep=4pt, font=\small, line width=0.7pt}]

  % ---- Top row: feature flow ----
  \node[box, fill=blue!10, minimum width=1.8cm] (fin)  at (0.0, 0) {$\mathbf{h}$};
  \node[op,  fill=orange!25, draw=orange!60]    (mul)  at (2.6, 0) {$\odot$};
  \node[op,  fill=orange!25, draw=orange!60]    (add)  at (4.6, 0) {$+$};
  \node[box, fill=blue!10, minimum width=1.8cm] (fout) at (7.0, 0) {$\mathbf{h}'$};

  % ---- γ and β directly below mul and add ----
  \node[box, fill=orange!18, draw=orange!50, minimum width=1.5cm]
    (gam) at (2.6, -1.4) {$\boldsymbol{\gamma}{\in}\mathbb{R}^C$};
  \node[box, fill=orange!18, draw=orange!50, minimum width=1.5cm]
    (bet) at (4.6, -1.4) {$\boldsymbol{\beta}{\in}\mathbb{R}^C$};

  % ---- MLP below centre ----
  \node[box, fill=orange!10, draw=orange!40, minimum width=3.0cm]
    (mlp) at (3.6, -2.7) {MLP\enspace{\scriptsize $\mathbb{R}^1 \!\to\! \mathbb{R}^{64} \!\to\! \mathbb{R}^{2C}$}};

  % ---- Rate scalar ----
  \node[box, fill=gray!8, minimum width=2.4cm]
    (r) at (3.6, -3.9) {$r = \log_2(\mathrm{bpp})$};

  % ---- Feature flow arrows ----
  \draw[->] (fin) -- (mul);
  \draw[->] (mul) -- (add);
  \draw[->] (add) -- (fout);

  % ---- Conditioning: straight vertical arrows ----
  \draw[->, orange!65, thick] (gam) -- (mul);
  \draw[->, orange!65, thick] (bet) -- (add);

  % ---- MLP splits to γ and β ----
  \draw[->] (mlp.north) -- ++(0,0.12) -| (gam.south);
  \draw[->] (mlp.north) -- ++(0,0.12) -| (bet.south);

  % ---- Rate to MLP ----
  \draw[->] (r) -- (mlp);

  % ---- Equation (above feature flow) ----
  \node[font=\small, above] at (3.4, 0.4)
    {$\mathbf{h}' = \boldsymbol{\gamma}(r) \odot \mathbf{h} + \boldsymbol{\beta}(r)$};
\end{tikzpicture}
\caption{%
  Feature-wise Linear Modulation (FiLM).
  A small MLP maps the conditioning signal $r$ (here: $\log_2(\mathrm{bpp})$)
  to per-channel scale $\boldsymbol{\gamma}$ and shift $\boldsymbol{\beta}$ vectors.
  These are applied channel-wise to the intermediate feature maps $\mathbf{h}$,
  conditioning the entire learned representation on the rate without modifying
  the backbone architecture.
}
\label{fig:film-diagram}
\end{figure}

\paragraph{Why FiLM over other conditioning approaches.}
Alternative conditioning strategies include:
\begin{itemize}
  \item \textit{Input concatenation}: append the rate scalar to each point's
    feature vector before processing.
    This only injects the conditioning signal at the first layer and requires
    the network to propagate it through multiple non-linear transformations.
  \item \textit{Output gating}: multiply the final displacement by a scalar
    function of the rate.
    This adjusts magnitude only, not the pattern of corrections.
  \item \textit{Separate models}: train one model per rate.
    This eliminates rate generalisation and increases parameter count linearly.
\end{itemize}
FiLM acts on intermediate features, conditioning the entire learned
representation on the rate. It adds minimal parameters (one MLP per modulated
layer) while providing maximum flexibility: the network can learn different
correction \emph{patterns} at different rates, not just different
magnitudes.

\subsection{Chamfer Distance}
\label{ssec:chamfer}

The Chamfer Distance (CD)~\cite{fan2017pointset} between two point sets $S$ and $T$ is:
\begin{equation}
  \CD(S, T) = \frac{1}{|S|}\sum_{s \in S} \min_{t \in T} \|s - t\|^2
            + \frac{1}{|T|}\sum_{t \in T} \min_{s \in S} \|t - s\|^2.
\end{equation}
The first term (forward direction, $S \to T$) penalises points in $S$ that
have no nearby counterpart in $T$: false negatives from $S$'s perspective.
The second term (reverse direction, $T \to S$) penalises points in $T$ with
no counterpart in $S$: false positives in $T$.

\paragraph{Differentiability.}
CD is differentiable with respect to the coordinates of points in $S$:
for each $s \in S$, its nearest neighbour $t^* = \arg\min_t \|s-t\|^2$ is
a constant (in the backward pass), so $\partial \|s - t^*\|^2 / \partial s = 2(s - t^*)$.
This makes CD usable as a training loss for point position regression.

\paragraph{Computational complexity.}
Naively, CD requires $O(|S||T|)$ nearest-neighbour queries.
For patches of $\sim$500 points, this is manageable; for full point clouds of
$\sim$800\,000 points, efficient KD-tree or GPU ball-query~\cite{qi2017pointnetplusplus} implementations
are required.
In this work, CD is computed on small patches during training; at evaluation
time, per-point D1 PSNR is used instead.

\paragraph{Relation to D1 PSNR.}
The D1 PSNR metric (Section~\ref{ssec:d1-psnr}) uses the maximum of the
two directional nearest-neighbour distances and normalises by the bounding
box diagonal, while CD uses the symmetric mean and does not normalise.
CD is more symmetric and better suited as a training loss; D1 PSNR is more
interpretable as a quality metric.


% ============================================================
